% !TeX root = ../papers/nc-p3.tex
\section{階数二以下の曲面成分と残余六次曲面}
\label{sec:low-rank-surface-factors}
\begin{theorem}\label{thm:low-rank-surface-factors}
四辺有限の標準有限Frobenius窓で入力 $X$ は被約、全二次式を成分に持つ反対平方根 $N$ は一般可逆、
階数一軌跡は有限とする。階数二以下の軌跡の既約曲面成分は平面に限られる。
各成分平面 $h=0$ では出力二次式が二つの線型埋込みで与えられる別々の平面の点の積になる。
その平面上の印は高々二つで、残余六次式 $\det N/h^2$ は平面外の少なくとも十八の印を三重に通る。
この六次式が既約かつ被約ならその曲面は非正規である。
\end{theorem}
\begin{proof}[因子の次数]
代数閉包上で既約成分の式を $f$、その次数を $e$ とする。$(f)$ で局所化して単元行列を消すと
$f^2\mid\det N$、従って $e\le4$。各印では同様に行列式の位数は少なくとも3である。
$e=4$ なら全ての印で $f$ が特異になる。その偏導三次式は入力三次評価の単射性で全て零となり矛盾する。
$e=3$ なら残余二次式の外には少なくとも六つの印がある。これは
定理\ref{thm:reducible-octic-factors}の証明中の二次評価と隣辺の完全跡による上界で、低階数軌跡の有限性を使わない。
これらは $f$ の特異点である。既約三次の孤立特異点は高々四つ、特異曲線があれば特異集合は
一本の直線でそこに載る印は高々三つである（同定理の割線による証明）。ゆえに $e\le2$。
\end{proof}
\begin{proof}[二つの一次因子の順序と印の個数]
$S=V(f)$ はこの段階で正規である。二次式の因数に順序をつける射
$\mathbb P^3\times\mathbb P^3\to\mathbb P^9$ は階数二以下の像へ有限で、
平方の外で次数二のエタール被覆である（因子交換による商）。
$S$ の滑らかな部分から階数一の有限集合を除いた開集合は幾何学的に単連結である。
実際平面・滑らかな二次曲面は単連結、二次錐の頂点外は射影直線上のアフィン直線束であり、
滑らかな複素曲面から有限点を除いても基本群は変わらない。標数零では係数を有限生成部分体に降ろし
複素数体で判定できる。よって順序のついた各成分は $S$ に有限双有理で、その被約閉包は正規性により
$S$ に同型である。二因子の射の超平面束を $A,B$ とすれば、大域生成され $A+B=2H$ となる。

$e=2$ とする。制限した行列成分の線型包は九次元である。因子像のベクトル包の次元を $r,s$ とすると、
片方が2以下なら積の包は高々8次元である。両方が3でも包の交わりは二次元以上なので、
その外積が二つの包のテンソル積から対称積への核に入り、やはり高々8次元である。
従って $r,s\ge3$ で一方は4である。滑らかな二次曲面で二次数 $(i,j),(2-i,2-j)$ を比べると
両者とも $(1,1)$ となる。実際 $(2,0),(0,2)$ は両方とも三切断、他の分割は一方が二切断以下である。
錐でもPicard群は $H$ の生成する無限巡回群である（$\mathbb F_2$ 上で縮約曲線に直交する部分を取る）。
従って両方の射は周囲の線型写像で、一方は可逆である。両像が一致する点はその相対線型変換の
固有空間の射影と $S$ の交わりにある。核の点は後者の射の基点にはならない。
階数一軌跡が有限なので射影平面以上の固有空間はなく、固有直線も $S$ に含まれない。
ベクトル次元2の固有空間は高々二印、次元1なら高々一印を含み、相異なる固有空間の直和の次元は4以下である。
従って $S$ 上の印は高々四つ、残余四次式が少なくとも十六印を三重に通る。

そのような四次式は存在しない。平面因子を持てばその上の印は七つ以下、外の三重点は残余三次の全二階偏導の
共通零点にある。その一次式の少なくとも一つは非零なのでこれも七つ以下である。
平面因子がなく二つの既約二次式に分かれる場合、三重点は少なくとも一方の頂点である（重複因子でも同じ）。
既約四次なら非共線な三つの三重点を含めない。三点の平面切断では各二点を結ぶ直線が少なくとも二重に出る。
一回割った後も次数3に対して両端の重複度の和が4以上あるからである。従って三直線で次数を超え、
平面因子が生じてしまう。これで二次因子も除いた。

$S$ が平面なら積の包は六次元で、一方が定数の分割では高々四次元である。
従って $A=B=\mathcal O(1)$ で両方は線型埋込みである。像平面が一致すればその共通零化列 $v$ について
$Nv$ は $h$ で割れ、定理\ref{thm:opposite-column-divisor}に反する。
異なる像平面の交線への二つの逆像直線上でのみ両点は一致し得る。逆像直線が一致する場合も
直線の二つの射影同型の一致点は高々二つか直線全体で、後者は階数一の有限性に反する。
従って平面外に少なくとも十八の印が残る。
\end{proof}
\begin{proof}[残余六次の正規性の排除]
正規な整六次曲面 $G$ が十八以上の三重点を持つと仮定し、解消を $Y$ とする。
正規Gorenstein有理曲面特異点は有理二重点で重複度2であるから、各三重点は局所幾何種数が1以上である。
従って $\chi(\mathcal O_Y)\le\chi(\mathcal O_G)-18=11-18=-7$。
標数零の滑らかな射影曲面の分類より $Y$ は種数8以上の曲線上の線織曲面に双有理である。
一般平面切断は滑らかな種数10の六次曲線で、底への射の次数 $b$ は正である。
実際この次数は一般繊維の入力超平面次数である。Riemann--Hurwitzより $18\ge b(2\cdot8-2)$、
従って $b=1$。一般繊維は入力直線になる。有限な特異点の一つを一般の直線が通れば $G$ は錐になり、
正規な錐の特異点はその頂点のみで矛盾する。そうでなければ一般直線は滑らかな部分にあり、
$K_G=2H$ と随伴公式より解消上の自乗が $-4$ となって、動く直線であることに反する。
\end{proof}

特に平面上の二因子を一次列 $u,v$ に延長すると $N=uv^t+vu^t+hC$ と書け、$C$ は線型対称行列である。
残余六次は符号を除き次の線型対称表示を持つ。
\[
\det\begin{pmatrix}C&u&v\\ u^t&0&-h\\ v^t&-h&0\end{pmatrix}
                 =-\det N/h^2 .
\]
これは $h\ne0$ で右下二行を消去すれば得られる多項式恒等式である。
